\section{Introduction}
\label{sec:introduction}

The estimation of a time-varying conditional quantile from a single observed time series is a fundamental problem in nonparametric statistics, arising in environmental monitoring, biomedical signal processing, finance, and astrophysics, among other fields. Unlike mean regression, quantile estimation provides a direct view of the upper or lower tails of the response distribution, which are often of greater scientific interest than the central tendency. When the underlying trend is believed to be piecewise polynomial---constant level shifts, piecewise linear behavior, or piecewise quadratic curvature---trend filtering provides a flexible and computationally tractable framework, producing estimates that adapt automatically to abrupt changes while remaining smooth in regions of slow variation.

The frequentist approach to quantile trend filtering was developed by \citet{brantley2020baseline}, who combined the check loss of \citet{koenker1978regression} with the discrete difference penalty of \citet{tibshirani2014adaptive} and solved the resulting nonsmooth convex program via the alternating direction method of multipliers (ADMM). Their estimator scales well to large signals and inherits the locally adaptive properties of $\ell_1$ trend filtering. Theoretical properties, including risk bounds, have been studied by \citet{madrid2022risk}. This approach provides a natural benchmark for estimation, but its output is a point estimate at a selected penalty level. Regularization selection, for example through an extended Bayesian information criterion, and uncertainty quantification require procedures beyond the ADMM optimization itself.

We develop an adaptive, fully Bayesian extension of quantile trend filtering computed by ADMM. The connection is exact: under an asymmetric Laplace likelihood and a global Laplace prior on the trend differences, the conditional posterior mode at $\sigma=1$ and fixed $\lambda$ coincides with the frequentist estimator. We extend this formulation by learning the regularization strength and likelihood scale jointly with the trend, and use ADMM within posterior sampling to retain its useful computational structure. This retains a posterior distribution over the trend while avoiding the growing collection of latent-scale updates and repeated matrix computations that make conventional augmented Gibbs implementations costly for long time series. Section~\ref{subsec:individual} establishes the mode equivalence.

Related progress on Bayesian quantile regression in additive and spline settings has been made by \citet{fasiolo2021fast}, though their approach targets smooth additive models rather than piecewise polynomial trends of arbitrary order. The analogous trajectory for mean trend filtering has been traced by Heng, Zhou, and Chi, who placed the $\ell_1$ trend filtering problem of \citet{tibshirani2014adaptive} in a Bayesian framework using proximal Markov chain Monte Carlo techniques \citep{pereyra2016proximal}. The quantile setting introduces an additional computational difficulty: both the check loss and the difference penalty are nonsmooth. The challenge is therefore to preserve the useful optimization structure while making joint posterior inference computationally practical.

The Bayesian quantile trend filtering method of \citet{onizuka2024fast} illustrates this computational challenge. Its Gibbs sampler uses latent-variable representations of the likelihood and shrinkage prior, requiring observation-level latent updates, local-scale updates for the trend differences, and a Gaussian trend update with a precision matrix that changes each iteration. These updates support flexible local shrinkage, especially under the horseshoe prior, but impose substantial computational overhead. In the supplied implementation, dense inversion and Gaussian simulation remain a bottleneck even with the tested sparse option enabled. Our formulation instead uses one global regularization parameter per quantile and an ADMM proximal system that permits factorization reuse.

The same authors also propose a faster mean-field variational Bayes (VB) approximation, together with a calibration procedure for its uncertainty summaries \citep{onizuka2024fast}. In exploratory runs beyond the $n=100$ reference setting, we encountered numerical failures in the public VB implementation. This implementation-specific observation provides additional motivation for a sampling alternative; our comparisons focus on estimation accuracy and computational performance rather than a quantitative assessment of uncertainty calibration.
% Writing note: document the VB repository commit, tested sample sizes, seeds, settings, and exact failures in the dedicated computational section before presenting this exploratory observation as a benchmark result.

The present paper combines the mode connection with a computational use of ADMM: the same class of $\ell_1$ trend filtering subproblems becomes the engine of posterior simulation. We construct a three-block Gibbs scheme that updates the single global regularization parameter $\lambda$ and the likelihood scale $\sigma$ from conjugate full conditionals, and updates the trend through a proximal Metropolis-adjusted Langevin (P-MALA) step \citep{pereyra2016proximal}. This formulation retains the local adaptation induced by the $\ell_1$ difference penalty without introducing a separate random shrinkage scale for every difference or a latent variable for every observation. For each quantile, only two scalar hyperparameters are sampled in addition to the trend. The proximal subproblems are solved using ADMM. In the constrained formulation, their coefficient matrix depends only on the difference operator and fixed algorithmic parameters, permitting reuse of a banded factorization. With sparse operations and a bounded number of inner ADMM iterations, the work per outer iteration is linear in the sample size. The advantage is thus a simpler hierarchy and reusable optimization structure, together with posterior inference for the trend and its regularization strength.

The framework extends naturally to joint inference across multiple quantile levels under the standard non-crossing constraint. The single-quantile P-MALA step is replaced by a sequential update in which the proximal map is solved under box constraints determined by the most recent values of the neighboring quantile trends. The required modification to the ADMM solver introduces a single auxiliary variable and an elementwise projection step, leaving its overall structure and computational profile intact. This joint construction, however, also exposes a difficulty intrinsic to gradient-based samplers under the asymmetric Laplace likelihood at extreme quantile levels. For $\tau$ close to zero or one, nearly all residuals share the same sign and the gradient of the smoothed check loss approaches a constant vector, so the P-MALA proposal carries negligible directional information and the chain develops a persistent drift away from the data. The locally adaptive prior of \citet{onizuka2024fast} partially mitigates this pathology, and their sampler produces degraded but not catastrophic estimates at extreme quantile levels. We instead resolve it through a strategy that exploits the box-constraint structure of the joint sampler: we augment the set of inferential quantile levels with auxiliary \emph{sentinel} quantiles placed beyond the target range, whose role is purely structural. The sentinels are sampled as ordinary members of the joint posterior and provide a dynamically updating upper or lower bound on the extreme target quantile through the non-crossing constraint, confining its exploration to a statistically sensible region at every iteration. The mechanism is computationally cheap and, in our simulation study, allows the proposed sampler to match or slightly improve upon the extreme-quantile accuracy of \citet{onizuka2024fast} while using a globally penalized prior that is structurally simpler than the per-coefficient horseshoe.

The main contributions of this work are as follows.
\begin{enumerate}
    \item We develop an adaptive Bayesian extension of the frequentist quantile trend filtering estimator: its conditional posterior mode recovers the ADMM solution at $\sigma=1$ and fixed $\lambda$, while the full hierarchy learns both parameters jointly with the trend and propagates their uncertainty into posterior summaries.
    \item We use ADMM within proximal posterior sampling, with one global regularization parameter per quantile and no observation-level likelihood augmentation. The constrained proximal system permits reuse of a banded factorization; sparse operations give linear work in $n$ for a fixed inner iteration budget, avoiding the dense inverse and Gaussian simulation used by the supplied Gibbs benchmark.
    \item We extend the sampler to the joint estimation of multiple quantile levels under non-crossing constraints, incorporating the constraints exactly through a box-projection step inside the inner ADMM solver rather than through approximate or penalty-based enforcement.
    \item We introduce the sentinel quantile strategy, a structurally minimal device that resolves the gradient degeneracy of proximal MALA at extreme quantile levels and recovers competitive accuracy at the tails of the response distribution. The mechanism is generic and applies to any gradient-based sampler for quantile trend filtering, not only the one developed here.
    \item Through a comprehensive simulation study spanning four signal structures, three noise distributions, three difference orders, and both interior and extreme quantile levels, we demonstrate that the proposed sampler delivers accuracy comparable to or better than the competitive Bayesian approach of \citet{onizuka2024fast}, with posterior uncertainty summaries and reduced computational overhead relative to the supplied dense Gibbs implementation, and that it does so using a single global penalty parameter rather than the per-coordinate local shrinkage employed by competing priors, which simplifies the prior specification and the resulting sampler.
\end{enumerate}

The remainder of the paper is organized as follows. Section~\ref{sec:methodology} develops the methodology, presenting the data model and likelihood, the individual-quantile sampler with its three conjugate and proximal updates, the joint sampler with non-crossing constraints and the associated ADMM proximal operator, and the sentinel quantile strategy. Section~\ref{sec:simulations} reports the simulation study comparing the proposed sampler with the relevant frequentist and Bayesian competitors across interior and extreme quantile levels. Section~\ref{sec:discussion} discusses the results, limitations, and avenues for extension.
